\exercisesection{§ 2 的习题}

记 R 为实向量空间 V 中的约化根系。

\begin{enumerate}
\def\labelenumi{\arabic{enumi})}
\tightlist
\item
  令 \(x \in V\)，令 \(W(x)\) 为 \(W(R)\) 中满足下述条件的元素 \(w\) 所成的子群：
\end{enumerate}

\[
w (x) - x \in \mathrm{Q} (\mathrm{R}).
\]

证明 \(W(x)\) 由反射生成。（利用 \(R^{\prime}\) 的仿射 Weyl 群。）

\begin{enumerate}
\def\labelenumi{\arabic{enumi})}
\setcounter{enumi}{1}
\item
  假设 R 不可约。令 \(\{\alpha_{1},\ldots,\alpha_{l}\}\) 为 R 的一个基，令 \(\tilde{\alpha}=\sum_{i}n_{i}\alpha_{i}\) 为 R 的最高根，令 f 为 R 的连通指标。证明 f-1 等于满足 \(n_{i}=1\) 的指标 i 的数目。
\item
  在第 3 号的记号和假设下，令 u 为仿射空间 E 的一个置换超平面 \(L_{\alpha,k}\)（\(\alpha \in R, k \in Z\)）的自同构。证明 u 是一个平移。（若 \(u_{0}\) 是与 u 关联的线性映射，证明 u 的转置使 R 稳定，因而属于群 A(R)。）由此推出 G 是仿射空间 E 的自同构群中 \(W_{a}\) 的正规化子。
\item
  令 \(C'\) 为相对于 \(W(R)\) 的 \(V^{*}\) 中的一个室，令 C 为包含于 \(C'\) 且顶点为 0 的胞腔。令 \(S_{a}\)（分别为 S）为 C（分别为 \(C'\)）的壁上的反射所成的集合。对 \((W_{a}, S_{a})\) 和 \((W, S)\) 而言，它们是 Coxeter 系统，且 \(S \subseteq S_{a}\)。证明 \(w \in W_{a}\) 是 \((S, \varnothing)\)-约化的（第 IV 章 §3，习题 3），当且仅当 \(w(C) \subseteq C'\)。
\end{enumerate}

¶ 5) 假设 R 不可约，并在 V 中选取 R 的一个室 C；以 B 表示 R 的相应基。

\begin{enumerate}
\def\labelenumi{\alph{enumi})}
\item
  证明 R 的极小权（§1，习题 24）在 P(R) 中构成 P(R)/Q(R) 的非零元素的一组代表元。（将命题 6 的推论应用于 R'。）
\item
  令 X 为 Q(R) 的一个饱和子集（§ 1，习题 23）。假设 X 非空且不只含 \{0\}。令 p 为 X 中长度最小的非零元素。证明 \(p \in R\)。（根据 a)，p 不满足 § 1 习题 24 的条件 (ii)；因此存在 \(\alpha \in R\)，使得 \(\langle p, \alpha^{-} \rangle \geqslant 2\)，于是 \(p - \alpha \in \mathbf{X}\)。由于 \(p - \alpha\) 的长度严格小于 \(p\) 的长度，\(p - \alpha = 0\)，所以 \(p \in \mathbb{R}\)。）
\item
  令 p 为不属于 Q(R) 的支配权。证明由 p 生成的饱和子集 S(p)（参见 §1，习题 23）包含唯一的极小权。（注意 S(p) 包含于模 Q(R) 的一个非平凡类中；应用 §1 的 a) 和习题 24 c) 得出结论。）
\item
  令 p 为不属于 Q(R) 的支配权。证明下列两个性质等价：
\end{enumerate}

\begin{enumerate}
\def\labelenumi{(\roman{enumi})}
\item
  \(p\) 是极小权。
\item
  不存在支配权 \(q \neq p\)，使得 p - q 是 B 中元素的具有非负整数系数的线性组合。
\end{enumerate}

（若 \(q\) 满足 (v) 的条件，令 \(p_1\) 为属于 \(S(q)\) 的极小权。则 \(p_1 \equiv p \mod Q(R)\)。\(p_1 < p\)，而 \(a\)) 表明 \(p\) 不是极小权。因此，(i) \(\Longrightarrow\) (v)。反之，若 \(p\) 不是极小权，令 \(q \in S(p) - W.p\)；必要时用 W 的元素变换 \(q\)，可设 \(q \in \overline{C}\)；由 §1 的习题 23，\(q < p\)。\(q \neq p\)，且 \(q \equiv p \mod Q(R)\)。因此，(v) \(\Longrightarrow\) (i)。）

